\section{Contoh Terhitung}
\label{sec:contoh-rsa}

Seluruh perhitungan pada bagian ini memakai parameter yang sama, yaitu $p = 47$,
$q = 71$, dan $e = 79$, sehingga kunci publiknya $(e, n) = (79, 3337)$ dan kunci
privatnya $(d, n) = (1019, 3337)$. Modulusnya sengaja dipilih kecil supaya setiap langkah
dapat diperiksa dengan tangan, tetapi seluruh sifat yang diperlihatkan berlaku sama pada
modulus sepanjang 2048 bit. Keempat contoh berikut disusun berurutan: membangkitkan kunci,
mengenkripsi satu blok, mengenkripsi pesan yang lebih panjang daripada modulusnya, lalu
memperlihatkan dua kelemahan RSA telanjang yang menjelaskan mengapa \textit{padding}
diperlukan.

\begin{example}[Pembangkitan kunci RSA]
\label{ex:rsa-kunci}
Alice memilih dua bilangan prima kecil $p = 47$ dan $q = 71$.
\begin{enumerate}
    \item Hitung modulus:
    \[ n = p \cdot q = 47 \times 71 = 3337. \]
    \item Hitung fungsi Totient Euler:
    \[ \varphi(n) = (p-1)(q-1) = 46 \times 70 = 3220. \]
    \item Pilih kunci publik $e = 79$; syaratnya terpenuhi karena $1 < 79 < 3220$ dan
          $\text{PBB}(79, 3220) = 1$.
    \item Hitung kunci privat $d$ sebagai invers $e$ modulo $\varphi(n)$:
    \[ e \cdot d \equiv 1 \pmod{3220} \;\Longrightarrow\; d = \frac{1 + k\,\varphi(n)}{e}
       = \frac{1 + 3220k}{79}. \]
    Nilai $k = 1, 2, 3, \dots$ dicoba hingga pembilang habis dibagi 79; untuk $k = 25$,
    \[ d = \frac{1 + 3220 \times 25}{79} = \frac{80501}{79} = 1019. \]
    Pemeriksaan: $79 \times 1019 = 80501 = 1 + 25 \times 3220$.
\end{enumerate}
Kunci publik Alice adalah $(e, n) = (79, 3337)$ dan kunci privatnya $(d, n) = (1019, 3337)$;
hanya $n$ dan $e$ yang boleh disebarkan. Prosedur yang sama dijalankan dengan Algoritma
Euclidean diperluas pada Contoh~\ref{ex:euclid-rsa}, yang menghasilkan nilai $d$ yang sama
dalam lima pembagian saja, tanpa perlu menebak nilai $k$.
\end{example}

\begin{example}[Enkripsi dan dekripsi RSA]
Bob mengirim pesan $m = 16$ kepada Alice menggunakan kunci publik $(e, n) = (79, 3337)$.
\begin{enumerate}
    \item Enkripsi oleh Bob:
    \[ c = m^e \bmod n = 16^{79} \bmod 3337 = 1493. \]
    \item Dekripsi oleh Alice dengan kunci privat $(d, n) = (1019, 3337)$:
    \[ m = c^d \bmod n = 1493^{1019} \bmod 3337 = 16. \]
\end{enumerate}
Plainteks $m = 16$ pulih dengan tepat. Perhatikan bahwa bilangan $16$ di sini adalah
sebuah bilangan bulat, bukan huruf; cara mengubah teks menjadi bilangan bulat dibahas
pada contoh berikutnya.
\end{example}

\begin{example}[Enkripsi pesan multi-blok ``HELLO ALICE'']
\label{ex:rsa-hello-alice}
Bob ingin mengirim pesan $M = \text{HELLO ALICE}$ kepada Alice. Pesan itu harus diubah
lebih dahulu menjadi bilangan bulat, lalu dipecah menjadi blok-blok yang lebih kecil
daripada modulus.

\textbf{Langkah 1: pengodean.} Setiap huruf dikodekan menjadi dua digit menurut aturan
$\text{A} = 00$, $\text{B} = 01$, sampai $\text{Z} = 25$. Spasi diabaikan, sebagaimana
lazim pada contoh ini, dan dikembalikan kemudian ketika pesan diterjemahkan.
Tabel~\ref{tab:pengodean-hello} memperlihatkan hasilnya.

\begin{table}[htbp]
    \centering
    \begin{tabular}{@{}l *{10}{c}@{}}
    \toprule
    Huruf & H & E & L & L & O & A & L & I & C & E \\
    \midrule
    Kode & 07 & 04 & 11 & 11 & 14 & 00 & 11 & 08 & 02 & 04 \\
    \bottomrule
    \end{tabular}
    \caption{Pengodean pesan ``HELLO ALICE'' dengan aturan $\text{A}=00$ sampai
    $\text{Z}=25$. Spasi tidak dikodekan}
    \label{tab:pengodean-hello}
\end{table}

Deretan kodenya menjadi $m = 07041111140011080204$, yaitu bilangan 20 digit. Nilai itu
jauh lebih besar daripada $n = 3337$, sehingga harus dipecah. Dengan mengambil empat digit
per blok diperoleh $m_1 = 0704$, $m_2 = 1111$, $m_3 = 1400$, $m_4 = 1108$, dan
$m_5 = 0204$. Kelimanya memenuhi syarat $0 \le m_i \le 3336$ pada
Persamaan~\eqref{eq:syarat-blok}.

\textbf{Langkah 2: enkripsi.} Setiap blok dienkripsi dengan kunci publik Alice,
$c_i = m_i^{79} \bmod 3337$. Hasilnya adalah

\[
\begin{aligned}
m_1 = 0704 &\quad\longrightarrow\quad c_1 = 704^{79} \bmod 3337 = 0328, \\
m_2 = 1111 &\quad\longrightarrow\quad c_2 = 1111^{79} \bmod 3337 = 0301, \\
m_3 = 1400 &\quad\longrightarrow\quad c_3 = 1400^{79} \bmod 3337 = 2653, \\
m_4 = 1108 &\quad\longrightarrow\quad c_4 = 1108^{79} \bmod 3337 = 2986, \\
m_5 = 0204 &\quad\longrightarrow\quad c_5 = 204^{79} \bmod 3337 = 1164 .
\end{aligned}
\]

Bob mengirimkan cipherteks $C = 0328, 0301, 2653, 2986, 1164$ kepada Alice.

Perhatikan penulisan $c_2 = 0301$ dan bukan $301$. Nilai $c_2$ memang hanya tiga digit,
tetapi ia harus tetap dituliskan dalam empat digit agar sejajar dengan blok-blok yang
lain. Bila nol di depan itu dibuang, Alice -- atau program apa pun yang membaca
cipherteks itu sebagai satu deretan panjang -- akan kehilangan batas antrablok dan
memperoleh hasil yang kacau. Aturan yang sama berlaku pada pengodean huruf A, yang
nilainya $00$ dan bukan sekadar kosong.

\textbf{Langkah 3: dekripsi.} Alice memakai kunci privatnya, $m_i = c_i^{1019} \bmod 3337$,
dan memperoleh kembali kelima blok semula. Tabel~\ref{tab:jejak-hello-alice} merangkum
seluruh perjalanan pesan dari huruf sampai huruf kembali.

\begin{table}[htbp]
    \centering
    \begin{tabular}{@{}c c c c l@{}}
    \toprule
    \textbf{Blok} & \textbf{Plainteks $m_i$} & \textbf{Cipherteks $c_i$} & \textbf{Pulih} &
    \textbf{Huruf} \\
    \midrule
    1 & \code{0704} & \code{0328} & \code{0704} & H, E \\
    2 & \code{1111} & \code{0301} & \code{1111} & L, L \\
    3 & \code{1400} & \code{2653} & \code{1400} & O, A \\
    4 & \code{1108} & \code{2986} & \code{1108} & L, I \\
    5 & \code{0204} & \code{1164} & \code{0204} & C, E \\
    \bottomrule
    \end{tabular}
    \caption{Jejak enkripsi dan dekripsi pesan ``HELLO ALICE'' dengan kunci publik
    $(79, 3337)$ dan kunci privat $(1019, 3337)$}
    \label{tab:jejak-hello-alice}
\end{table}

Deretan angka yang dipulihkan adalah $07041111140011080204$, yang diterjemahkan kembali
menjadi $\text{HELLOALICE}$. Spasi yang tadi diabaikan kini harus disisipkan kembali oleh
Alice berdasarkan pengetahuannya tentang bahasa, sebab tidak ada satu pun bit di dalam
cipherteks yang menyimpan informasi tentang posisi spasi.
\end{example}

\begin{example}[Mengapa RSA telanjang berbahaya]
\label{ex:rsa-telanjang}
Dua sifat RSA tanpa \textit{padding} dapat diperlihatkan langsung dengan parameter yang
sama, dan keduanya menerangkan mengapa skema \textit{padding} pada
Subbagian~\ref{subsec:padding-oaep} bukan pilihan melainkan keharusan.

\textbf{Sifat pertama: deterministik.} Karena rumus $c = m^e \bmod n$ tidak memuat unsur
keacakan sama sekali, plainteks yang sama selalu menghasilkan cipherteks yang sama. Pesan
``CDCD'' dikodekan menjadi dua blok kembar $m = 0203$; keduanya menghasilkan cipherteks
$c = 0203^{79} \bmod 3337 = 1488$. Penyerang yang menyadap deretan cipherteks itu
langsung mengetahui bahwa dua blok pertama pesan tersebut identik, tanpa perlu membuka
apa pun.

Lebih jauh, karena $e$ dan $n$ bersifat umum, penyerang dapat menyusun kamus terjemahan
seluruh kandidat pesan yang mungkin. Tabel~\ref{tab:kamus-rsa} memperlihatkan kamus untuk
seluruh huruf abjad, yang dihitung hanya dari kunci publik Alice.

\begin{table}[htbp]
    \centering
    \footnotesize
    \begin{tabular}{@{}l r@{\ }r@{\ }r@{\ }r@{\ }r@{\ }r@{\ }r@{\ }r@{\ }r@{\ }r@{\ }r@{\ }r@{\ }r@{}}
    \toprule
    Huruf & A & B & C & D & E & F & G & H & I & J & K & L & M \\
    Kode  & 00 & 01 & 02 & 03 & 04 & 05 & 06 & 07 & 08 & 09 & 10 & 11 & 12 \\
    $c$ & 0000 & 0001 & 3139 & 0158 & 2497 & 0270 & 2086 & 1254 & 2807 & 1605 & 3269 & 2120 & 0760 \\
    \midrule
    Huruf & N & O & P & Q & R & S & T & U & V & W & X & Y & Z \\
    Kode  & 13 & 14 & 15 & 16 & 17 & 18 & 19 & 20 & 21 & 22 & 23 & 24 & 25 \\
    $c$ & 2066 & 1983 & 2616 & 1493 & 0827 & 2562 & 1265 & 0116 & 1249 & 0702 & 2524 & 3022 & 2823 \\
    \bottomrule
    \end{tabular}
    \caption{Kamus sandi untuk seluruh huruf abjad, dihitung hanya dari kunci publik
    $(79, 3337)$. Tabel semacam ini dapat dibangkitkan penyerang dalam waktu sekejap}
    \label{tab:kamus-rsa}
\end{table}

\textbf{Sifat kedua: kerapuhan terhadap perubahan.} Misalkan penyerang menyadap
cipherteks $c = 1493$, yang tidak ia ketahui adalah plainteks $m = 16$. Ia memilih
bilangan $s = 3$ dan menghitung $s^e \bmod n = 3^{79} \bmod 3337 = 158$. Kemudian ia
mengganti cipherteks yang disadap dengan

\[
c' = c \cdot s^e \bmod n = 1493 \times 158 \bmod 3337 = 2304 .
\]

Bila cipherteks palsu itu diteruskan kepada Alice, dekripsinya adalah
$2304^{1019} \bmod 3337 = 48$, sedangkan $m \cdot s = 16 \times 3 = 48$. Jadi penyerang
berhasil mengubah pesan Alice menjadi bilangan 48 tanpa pernah mengetahui bahwa pesan
aslinya adalah 16. Bila pesan itu menyatakan nomor rekening atau nominal transfer,
perubahan tersebut dapat berakibat sangat serius.

Kedua sifat itulah yang dihilangkan oleh OAEP. Dengan menyisipkan keacakan sebelum
pemangkatan, OAEP membuat enkripsi menjadi probabilistik sehingga kamus pada
Tabel~\ref{tab:kamus-rsa} tidak lagi berguna, dan membuat perubahan cipherteks tidak lagi
menghasilkan perubahan plainteks yang dapat diprediksi.
\end{example}
